% generated by roam.py handbook export -- не редактировать, править заметки
% 20261025215715-критерии_денниса_море.org (body 884785f28138)

\begin{theorem}[{Критерий Денниса–Море}]\label[theorem]{note:D65FB534-B7F1-4194-8E4B-D4ACAA89E6D2}
  Пусть

  \begin{itemize}
    \item $f \in C^{2,2}_M$~--- \hyperref[note:7E287F4B-AC3E-4CF5-9D32-3BC2BC52D01A]{липшицев гессиан}
    \item $\nabla f(x^{*}) = 0$, $\nabla^2 f(x^{*}) > 0$
    \item $x_{k+1} = x_k -B_k^{-1}\nabla f(x_k) = x_k + d_k$:  \hyperref[note:6BA18878-DE34-46A5-ABE2-EDA09CAE56A7]{квазиньютоновские методы}
    \item $x_k \rightarrow x^{*}$
  \end{itemize}

  Тогда $x_k$ сходится \hyperref[note:9FA20DA5-3A95-4683-9CB5-C6C9A903C86A]{сверхлинейно} $\Leftrightarrow \lim \frac{\|(B_k - \nabla^2 f(x_k))d_k\|}{\|d_k\|} = 0$
\end{theorem}

\begin{remark}
  Условие требует, чтобы $B_k$ приближала гессиан только вдоль направления шага $d_k$, а не целиком. Поэтому квазиньютоновские методы могут сходиться сверхлинейно, даже если $B_k \not\rightarrow \nabla^2 f(x^{*})$.
\end{remark}

\begin{proof}[Доказательство (только $\Leftarrow$)]
  Так как $f \in C^{2,2}_M$ и $\nabla^2 f(x^{*}) > 0$, то по \hyperref[note:CDA139F9-5438-41E1-96B1-11A751BE40C5]{локальной сходимости метода Ньютона}:
  \[
  \exists R,C\,:\ \forall x_0\,:\  \|x_0-x^{*}\| < R\ x_k \rightarrow x^*,\text{ причем }
  \|\tilde{x}_{k+1} - x^{*}\| \le C\|x_k - x^{*}\|^{2}
  \]
  Где $\tilde{x}_{k+1}$ получен \hyperref[note:C0ED5501-CFC6-42BA-A53A-1B346CBA545D]{методом Ньютона} за один шаг из $x_k$. Так как $x_k \rightarrow x^{*}$, при больших $k$ имеем $\|x_k - x^{*}\| < R$, и оценка применима к $x_0 := x_k$
  \[
  \tilde{x}_{k+1} = x_k - \left( \nabla^2 f(x_k) \right)^{-1} \nabla f(x_k) = x_k + d_k^N
  \]
  За $d_{k}^N$  мы обозначили шаг Ньютона
  \[
  d_k^N := - \left( \nabla^2 f(x_k) \right)^{-1} \nabla f(x_k)
  \]
  Пусть  $d_k$~--- наш шаг
  \[
  d_k := - B^{-1}_k \nabla f(x_k) = x_{k+1} - x_k
  \]

  Сперва нам нужно доказать, что
  \[
  d_k - d_k^N = o(\|d_k\|)
  \]

  Рассмотрим эту разность
  \[
  d_k - d_k^N = d_k + \left( \nabla^2 f(x_k)\right)^{-1} \nabla f(x_k)
  = \left( \nabla^2 f(x_k)\right)^{-1} \left( \nabla^2 f(x_k) d_k + \nabla f(x_k) \right)
  \]

  Тут мы воспользовались тем, что
  \[
  d_k = \left( \nabla^2f(x_k) \right)^{-1} \left( \nabla^2 f(x_k) \right) d_k
  \]
  Заметим, что
  \begin{multline*}
  x_{k+1} = x_k -B_k^{-1}\nabla f(x_k) \Leftrightarrow x_{k+1} - x_k = -B_k^{-1}\nabla f(x_k) \\
  \Leftrightarrow B_k(x_{k+1} -x_k ) = - \nabla f(x_k) \Leftrightarrow B_kd_k = - \nabla f(x_k)
  \end{multline*}
  Давайте подставим $-B_kd_k$ вместо $\nabla f(x_k)$
  \[
  \left( \nabla^2 f(x_k)\right)^{-1} \left( \nabla^2 f(x_k) d_k + \nabla f(x_k) \right)
  = \left( \nabla^2 f(x_k)\right)^{-1} \left( \nabla^2 f(x_k) d_k - B_kd_k \right)
  \]
  Итого
  \[
  d_k - d_k^N
  = \left( \nabla^2 f(x_k)\right)^{-1} \left( \nabla^2 f(x_k) d_k - B_kd_k \right)
  \]

  А значит и их нормы равны
  \begin{multline*}
  \| d_k - d_k^N\| =
  \left\| \left( \nabla^2 f(x_k)\right)^{-1} \left( \nabla^2 f(x_k) d_k - B_kd_k \right) \right\|
  \le\\ \left\|\left(\nabla^2 f(x_k)\right)^{-1} \right\|  \cdot
  \left\| \left( \nabla^2 f(x_k) - B_k \right) d_k\right\|
  \end{multline*}

  Заметим, что гессиан непрерывен, тогда
  \[x_k \rightarrow x^{*} \Rightarrow \nabla^2 f(x_k) \rightarrow \nabla^2 f(x^{*})\]
  А еще мы знаем, что $\lambda_{\min}$ тоже \hyperref[note:5A4CF475-D55C-4CAA-9A68-BD8D3B60AEA9]{непрерывная}, следовательно
  \[
  \lambda_{\min}\left(\nabla^2 f(x_k)\right) \rightarrow \lambda_{\min} \left(\nabla^2 f(x^*)\right) = \mu > 0
  \]
  И для достаточно больших $k$ верно, что
  \[
  \lambda_{\min}\left(\nabla^2 f(x_k)\right)  > \frac{\mu}{2}
  \]
  Тогда при этом условии
  \[
  \lambda_{\max}\left((\nabla^2 f(x_k))^{-1}\right) = \| (\nabla^2 f(x_k))^{-1} \| =
  \frac{1}{\lambda_{\min}\left(\nabla^2 f(x_k)\right)} < \frac{2}{\mu}
  \]
  Тогда
  \[
  0 \le \| d_k - d_k^N\| \le \frac{2}{\mu} \cdot  \left\| \left( \nabla^2 f(x_k) - B_k \right) d_k\right\|
  \Rightarrow \| d_k - d_k^N\| = O\left( \|(\nabla^2 f(x_k) - B_k)d_k\| \right)
  \]
  А мы знаем, что по условиям задачи
  \[
  \lim \frac{\|(B_k - \nabla^2f(x_k))d_k\|}{\|d_k\|} = 0
  \]
  Тогда
  \[
  \| d_k - d_k^N\| = o(\|d_k\|)
  \]

  Тогда
  \[
  \|x_{k+1} - x^{*}\| = \|x_k + d_k - x^{*}\| \le \|x_k + d_k^N - x^{*}\| + \|d_k - d_k^N\| =
  \|\tilde{x}_{k+1} - x^{*}\| + o(\|d_k\|)
  \]
  И по \hyperref[note:CDA139F9-5438-41E1-96B1-11A751BE40C5]{локальной квадратичной сходимости} метода Ньютона
  \[
  \|x_{k+1} - x^{*}\| \le
  \|\tilde{x}_{k+1} - x^{*}\| + o(\|d_k\|)
  \le C\|x_k - x^{*}\|^2 + o(\|d_k\|)
  \]

  Заметим, что $\|x_k - x^{*}\| \rightarrow 0$, тогда
  \[
  C\|x_k - x^{*}\|^2 + o(\|d_k\|) = o(\|x_k - x^{*}\|)  + o(\|d_k\|)
  \]

  \begin{remark}
    Что такое быть о малым
    \begin{multline*}
    \|x_{k+1} - x^{*}\| = o(\|x_k - x^{*}\|) \Leftrightarrow \\
    \Leftrightarrow \forall \varepsilon >0 \ \exists N:\ \forall n \ge N\ \|x_{n+1} - x^{*}\| \le \varepsilon \|x_n - x^{*}\| \le \varepsilon^{n+1-N} \|x_N - x^{*}\|
    \end{multline*}
  \end{remark}

  Осталось понять, что $\|d_k\| = O(\|x_k - x^{*}\|)$
  \begin{multline*}
  \|d_k\| = \|x_{k+1} - x_k\| = \|x_{k+1} - x^{*} + x^{*} - x_k\|
  \le \|x_{k+1} - x^{*}\|  + \|x^{*} - x_k\| \le  \\
  o(\|x_k - x^{*}\|) + \|x_k - x^{*}\| + o(\|d_k\|) =
  O(\|x_k - x^{*}\|) + o(\|d_k\|)
  \end{multline*}
  Для больших $k$ имеем $o(\|d_k\|) \le \frac{1}{2}\|d_k\|$, откуда $\frac{1}{2}\|d_k\| \le O(\|x_k - x^{*}\|)$, то есть $\|d_k\| = O(\|x_k - x^{*}\|)$.

  Теперь по предыдущим неравенствам докажем \hyperref[note:9FA20DA5-3A95-4683-9CB5-C6C9A903C86A]{сверхлинейность} \\
  Хотим
  \[
  \frac{\|x_{k+1} - x^{*}\|}{\|x_k - x^{*}\|} \rightarrow 0
  \]
  Оценим выражение сверху
  \[
  \frac{\|x_{k+1} - x^{*}\|}{\|x_k - x^{*}\|} \le
  \frac{C\|x_k - x^{*}\|^2 + o(\|d_k\|)}{\|x_k - x^{*}\|}
  = C\|x_k - x^{*}\| + \frac{o(\|d_k\|)}{\|x_k - x^{*}\|}
  \]
  Но мы поняли, что $\|d_k\| = O(\|x_k - x^{*}\|) \Rightarrow o(\|d_k\|) = o(\|x_k - x^{*}\|)$ \\
  Итого
  \[
  \frac{\|x_{k+1} - x^{*}\|}{\|x_k - x^{*}\|} \le
   C\|x_k - x^{*}\| + \frac{o(\|x_k - x^{*}\|)}{\|x_k - x^{*}\|} \rightarrow 0
  \]

  Каждое слагаемое стремится к 0; И по \hyperref[note:C0408924-F658-43DF-8794-0C9AA85E7580]{Тест корней} это сверхлинейная сходимость
\end{proof}
